% ===========================================================================
%  Reporte final — ¿En qué minuto se perdió la partida?
%  Héctor Manuel Morán Maldonado
%  Compilar (desde la carpeta del proyecto):
%     python LoL_Reporte/scripts/cifras.py      (cifras y tablas desde los resultados)
%     latexmk -cd -lualatex LoL_Reporte/reporte_final.tex
%  Requiere LuaLaTeX (TeX Live o MiKTeX) y la fuente DejaVu Sans Mono.
% ===========================================================================
\documentclass[11pt,letterpaper]{article}

\usepackage{fontspec}
\usepackage[spanish,es-tabla,es-nodecimaldot]{babel}
\usepackage{microtype}
\usepackage[margin=2.5cm,top=2.6cm,bottom=2.6cm]{geometry}
\usepackage{graphicx}
\usepackage{booktabs}
\usepackage{array}
\usepackage{longtable}
\usepackage{float}
\usepackage{xcolor}
\usepackage{amsmath}
\usepackage{enumitem}
\usepackage{needspace}
\usepackage{caption}
\usepackage{fancyhdr}
\usepackage{titlesec}
\usepackage{fvextra}
\usepackage[hidelinks]{hyperref}

\newfontfamily\codigo{DejaVu Sans Mono}[Scale=0.78]
\graphicspath{{../LoL_Analisis/figuras/}}
\emergencystretch=2em

\definecolor{azul}{HTML}{2A78D6}
\definecolor{naranja}{HTML}{EB6834}
\definecolor{tinta}{HTML}{1A1A19}
\definecolor{suave}{HTML}{52514E}
\definecolor{papel}{HTML}{F4F6F8}

\captionsetup{font=small,labelfont={bf,color=azul},width=.92\linewidth}
\titleformat{\section}{\Large\bfseries\color{tinta}}{\thesection}{.7em}{}
\titleformat{\subsection}{\large\bfseries\color{suave}}{\thesubsection}{.7em}{}
\titlespacing*{\section}{0pt}{2.2ex plus 1ex minus .2ex}{1.2ex plus .2ex}

\pagestyle{fancy}\fancyhf{}
\renewcommand{\headrulewidth}{.4pt}
\fancyhead[L]{\footnotesize\color{suave}Reporte final · ¿En qué minuto se perdió la partida?}
\fancyhead[R]{\footnotesize\color{suave}H. M. Morán Maldonado}
\fancyfoot[C]{\footnotesize\color{suave}\thepage}

\newsavebox{\cajaresp}
\newenvironment{respuesta}[1]
  {\par\medskip\setlength{\fboxsep}{10pt}%
   \begin{lrbox}{\cajaresp}%
   \begin{minipage}{\dimexpr\linewidth-2\fboxsep-3pt\relax}%
   \color{tinta}\textcolor{azul}{\textbf{#1}}\par\smallskip\ignorespaces}
  {\end{minipage}%
   \end{lrbox}%
   \noindent\textcolor{azul}{\rule[-\dp\cajaresp]{3pt}{\dimexpr\ht\cajaresp+\dp\cajaresp\relax}}%
   \colorbox{papel}{\usebox{\cajaresp}}\par\medskip}

\newcommand{\dato}[1]{\textbf{\textcolor{azul}{#1}}}
\newcommand{\figura}[3]{\begin{figure}[!htbp]\centering
  \includegraphics[width=#2\linewidth]{#1}\caption{#3}\label{fig:#1}\end{figure}}
\newcommand{\codigofuente}[2]{\subsection*{#1}\addcontentsline{toc}{subsection}{#1}
  \VerbatimInput[breaklines, breakanywhere, numbers=left, numbersep=6pt,
    fontsize=\scriptsize, formatcom=\codigo, frame=lines, framesep=4pt]{#2}}

\newcommand{\nCuentas}{6}
\newcommand{\listaTotal}{2\,442}
\newcommand{\listaCola}{1\,251}
\newcommand{\nUnicas}{1\,241}
\newcommand{\nFinal}{1\,195}
\newcommand{\nDerrotas}{537}
\newcommand{\nVictorias}{658}
\newcommand{\nQuitadas}{46}
\newcommand{\nDuo}{13}
\newcommand{\nCuentaUno}{102}
\newcommand{\derCuentaUno}{41}
\newcommand{\nBottom}{839}
\newcommand{\pctBottom}{70}
\newcommand{\nParches}{46}
\newcommand{\mesInicio}{2024-10}
\newcommand{\mesFin}{2026-10}
\newcommand{\nPuuid}{9\,852}
\newcommand{\nNombres}{10\,473}
\newcommand{\nFugas}{0}
\newcommand{\nEventos}{1\,366\,619}
\newcommand{\nCuadros}{382\,760}
\newcommand{\pctUnicosSinCampeon}{1{,}0}
\newcommand{\pctPartidaUnica}{97}
\newcommand{\nPartidasPublicas}{1\,241}
\newcommand{\nEventosPublicos}{29\,696}
\newcommand{\eUnoPct}{67}
\newcommand{\eUnoLo}{63}
\newcommand{\eUnoHi}{71}
\newcommand{\eUnoConDdos}{535}
\newcommand{\eUnoDesac}{360}
\newcommand{\eUnoMedDif}{6}
\newcommand{\dUnoMedDer}{5}
\newcommand{\dUnoMedVic}{5}
\newcommand{\dUnoAntesCincoDer}{46}
\newcommand{\dUnoAntesCincoVic}{41}
\newcommand{\dDosMed}{13}
\newcommand{\derrotasArriba}{11}
\newcommand{\derrotasNuncaDetras}{2}
\newcommand{\hUnoDerTasa}{0{,}54}
\newcommand{\hUnoVicTasa}{0{,}87}
\newcommand{\hUnoDerPartidas}{126}
\newcommand{\hUnoVicPartidas}{564}
\newcommand{\hUnoMenos}{38}
\newcommand{\hUnoDerA}{0{,}47}
\newcommand{\hUnoVicA}{0{,}54}
\newcommand{\hUnoDerB}{0{,}61}
\newcommand{\hUnoVicB}{0{,}95}
\newcommand{\hUnoDerC}{0{,}58}
\newcommand{\hUnoVicC}{1{,}10}
\newcommand{\hDosObsDer}{39{,}6}
\newcommand{\hDosRegDer}{45{,}6}
\newcommand{\hDosCorDer}{35{,}9}
\newcommand{\hDosPCorDer}{0{,}017}
\newcommand{\hDosPRegDer}{0{,}003}
\newcommand{\hDosObsVic}{28{,}2}
\newcommand{\hDosRegVic}{24{,}0}
\newcommand{\hDosCorVic}{19{,}0}
\newcommand{\hDosPCorVic}{0{,}000}
\newcommand{\hDosPRegVic}{0{,}993}
\newcommand{\hDosUnoObsDer}{26{,}4}
\newcommand{\hDosUnoRegDer}{45{,}9}
\newcommand{\hDosUnoCorDer}{25{,}0}
\newcommand{\hDosUnoPCorDer}{\(p = 0{,}177\)}
\newcommand{\hDosUnoPRegDer}{\(p < 0{,}001\)}
\newcommand{\hDosUnoObsVic}{18{,}8}
\newcommand{\hDosUnoRegVic}{22{,}9}
\newcommand{\hDosUnoCorVic}{13{,}1}
\newcommand{\hDosUnoPCorVic}{\(p < 0{,}001\)}
\newcommand{\hDosUnoPRegVic}{\(p = 0{,}004\)}
\newcommand{\alfaBonf}{0{,}0167}
\newcommand{\hDosPasaBonf}{sí}
\newcommand{\hDosAntesCinco}{163}
\newcommand{\hDosPrimerObj}{5{,}4}
\newcommand{\neutrales}{86}
\newcommand{\hTresCinco}{37}
\newcommand{\hTresDiez}{24}
\newcommand{\hTresQuince}{20}
\newcommand{\hTresVeinte}{12}
\newcommand{\hTresVeinticinco}{8}
\newcommand{\hTresTreinta}{4}
\newcommand{\hTresRiesgoTreinta}{154}
\newcommand{\hTresUmbral}{5}
\newcommand{\hTresPunto}{27}
\newcommand{\hTresPuntoP}{4{,}0}
\newcommand{\hTresPuntoLo}{2{,}1}
\newcommand{\hTresPuntoHi}{7{,}8}
\newcommand{\hTresPuntoRiesgo}{198}
\newcommand{\hTresPuntoRemontaron}{8}
\newcommand{\hTresPuntoDiez}{22}
\newcommand{\hTresPuntoDiezRiesgo}{205}
\newcommand{\hTresPuntoTres}{35}
\newcommand{\hTresPuntoTresRiesgo}{51}
\newcommand{\pctAUno}{19}
\newcommand{\nAUno}{103}
\newcommand{\loAUno}{16}
\newcommand{\hiAUno}{23}
\newcommand{\pctADos}{10}
\newcommand{\nADos}{55}
\newcommand{\loADos}{8}
\newcommand{\hiADos}{13}
\newcommand{\pctATres}{4}
\newcommand{\nATres}{23}
\newcommand{\loATres}{3}
\newcommand{\hiATres}{6}
\newcommand{\pctACuatro}{15}
\newcommand{\nACuatro}{78}
\newcommand{\loACuatro}{12}
\newcommand{\hiACuatro}{18}
\newcommand{\pctSinClas}{52}
\newcommand{\nSinClas}{278}
\newcommand{\pctClas}{48}
\newcommand{\loSinClas}{48}
\newcommand{\hiSinClas}{56}
\newcommand{\eSeisCuatro}{64}
\newcommand{\eSeisSeis}{83}
\newcommand{\eSeisCuatroPct}{12}
\newcommand{\eSeisSeisPct}{15}
\newcommand{\eSieteCorte}{2025-06}
\newcommand{\eSieteUno}{209}
\newcommand{\eSieteDos}{328}
\newcommand{\eSieteMedUno}{12}
\newcommand{\eSieteMedDos}{14}
\newcommand{\eSieteP}{0{,}71}
\newcommand{\eOchoP}{0{,}19}
\newcommand{\eOchoPAprox}{0{,}19}
\newcommand{\eOchoEspMin}{1{,}7}
\newcommand{\botPartidas}{839}
\newcommand{\botDerrotas}{366}
\newcommand{\botHDosP}{0{,}034}
\newcommand{\botHTresTreinta}{4}
\newcommand{\botHUnoDerA}{0{,}42}
\newcommand{\botHUnoVicA}{0{,}55}
\newcommand{\botHUnoDerB}{0{,}60}
\newcommand{\botHUnoVicB}{0{,}96}
\newcommand{\botHUnoDerC}{0{,}60}
\newcommand{\botHUnoVicC}{1{,}12}
\newcommand{\exTranquilo}{78}
\newcommand{\exNuncaDiez}{74}
\newcommand{\exDdos}{19}
\newcommand{\exDur}{33}
\newcommand{\exDurClas}{30}
\newcommand{\vUnoArq}{A2}
\newcommand{\vUnoN}{55}
\newcommand{\vUnoDur}{34{,}1}
\newcommand{\vUnoMax}{5{,}9}
\newcommand{\vUnoMaxOro}{3\,500}
\newcommand{\vUnoMaxMin}{16}
\newcommand{\vUnoDdos}{34}
\newcommand{\vUnoArribaHasta}{33}
\newcommand{\vUnoArriba}{31}
\newcommand{\vUnoDurMin}{34}
\newcommand{\vUnoAbajoAntes}{7}
\newcommand{\vUnoAntesRemate}{2{,}1}
\newcommand{\vUnoMinAntesRemate}{32}
\newcommand{\vUnoRemateRival}{6}
\newcommand{\vUnoRematePropio}{0}
\newcommand{\vUnoPerdido}{64}
\newcommand{\vUnoFinal}{-1{,}7}
\newcommand{\vUnoObjPropios}{11}
\newcommand{\vUnoObjRival}{14}
\newcommand{\drN}{1\,132}
\newcommand{\drPcaDos}{83{,}4}
\newcommand{\drPcUnoCorr}{0{,}994}
\newcommand{\drLda}{76{,}4}
\newcommand{\drVveinte}{76{,}8}
\newcommand{\drSilMin}{0{,}15}
\newcommand{\drSilMax}{0{,}17}
\newcommand{\eaDefQuince}{1\,776}
\newcommand{\eaBotDer}{13}
\newcommand{\eaMidDer}{31}
\newcommand{\eaBotVic}{33}
\newcommand{\eaAbajoN}{247}
\newcommand{\eaAbajoYo}{40}
\newcommand{\eaAbajoMid}{52}
\newcommand{\eaMuerteContra}{1{,}30}
\newcommand{\eaMuerteFavor}{0{,}71}
\newcommand{\eaMiMuerte}{1{,}33}
\newcommand{\eaBaronContra}{1{,}48}
\newcommand{\eaCaminoMax}{0{,}003}
\newcommand{\eaMenosCincoQuince}{6}
\newcommand{\eaMuerteDer}{2{,}78}
\newcommand{\eaMuerteVic}{1{,}58}
\newcommand{\eaCampoDer}{26}
\newcommand{\eaCampoVic}{55}
\newcommand{\pgPartidas}{1\,195}
\newcommand{\pgPrueba}{232}
\newcommand{\pgNumMarc}{8}
\newcommand{\pgNumOro}{9}
\newcommand{\pgAntesMarc}{0{,}548}
\newcommand{\pgElegMarc}{0{,}536}
\newcommand{\pgAntesOro}{0{,}531}
\newcommand{\pgElegOro}{0{,}527}
\newcommand{\pgSoloOro}{0{,}545}
\newcommand{\pgXpDif}{0{,}002}
\newcommand{\pgObjDifMax}{0{,}0005}
\newcommand{\pgCalMax}{1}
\newcommand{\pgPruebaAntesMarc}{0{,}581}
\newcommand{\pgPruebaMarc}{0{,}574}
\newcommand{\pgPruebaOro}{0{,}562}


\begin{document}

% ===========================================================================
\begin{titlepage}
\thispagestyle{empty}
\centering
\vspace*{0.6cm}
\includegraphics[width=6.4cm]{assets/logo-up.png}

\vspace{1.1cm}
{\large\color{suave}Especialidad en Ciencia de Datos\par}
\vspace{0.25cm}
{\large\color{suave}Visualización de datos y narración de historias\par}

\vspace{1.5cm}
{\color{azul}\rule{\linewidth}{1.6pt}}
\vspace{0.85cm}

{\Huge\bfseries ¿En qué minuto se perdió la partida?\par}
\vspace{0.45cm}
{\LARGE Anatomía de mis derrotas en clasificatoria\par}
\vspace{0.35cm}
{\large\color{suave} Reporte final del proyecto · datos de la API de Riot Games\par}

\vspace{0.85cm}
{\color{azul}\rule{\linewidth}{1.6pt}}

\vfill

\begin{minipage}{0.46\linewidth}
  \raggedright
  {\color{suave}\small ALUMNO}\par\smallskip
  {\large Héctor Manuel Morán Maldonado\par}
\end{minipage}\hfill
\begin{minipage}{0.46\linewidth}
  \raggedleft
  {\color{suave}\small PROFESOR}\par\smallskip
  {\large Dr.\ Rodolfo Rubén Álvarez González\par}
\end{minipage}

\vspace{1.1cm}
{\color{suave}\rule{\linewidth}{.4pt}}
\vspace{0.45cm}

{\color{suave}\small
\nFinal{} partidas · \nCuentas{} cuentas propias · \nDerrotas{} derrotas · \nEventos{} eventos\\
Todas las cifras de este documento las generan los scripts del anexo a partir de los datos\par}
\end{titlepage}

\tableofcontents
\newpage

% ===========================================================================
\section*{Resumen}
\addcontentsline{toc}{section}{Resumen}

Reconstruí \nFinal{} partidas clasificatorias de mis \nCuentas{} cuentas de League of
Legends minuto a minuto, con la línea temporal de la API de Riot, para buscar en qué
minuto se pierde una partida. Medí la ventaja de oro relativa \(v(t)\) y contrasté tres
hipótesis contra mis propias victorias. Con ventaja, mis equipos toman menos objetivos
por minuto al perder, pero solo \emph{a partir del minuto 15} (H1). Que el quiebre de
la partida coincida con un monstruo del rival no se puede afirmar: con la definición
registrada no se distingue del azar, con la principal apenas lo supera, y en las
victorias la coincidencia es mayor (H2, no concluyente). Sí hay un punto casi sin
regreso, aunque se llega a él poco a poco: remontar estando abajo pasa del
\hTresCinco\,\% en el minuto 5 al \hTresQuince\,\% en el 15, y desde el minuto
\hTresPunto{} lo logra menos de una de cada veinte partidas (H3). Las cuatro formas de
perder que propuse explican el \pctClas\,\% de las derrotas; el \pctSinClas\,\% restante
son, sobre todo, partidas parejas que se deciden tarde. Cada cambio respecto al plan
registrado antes de ver los datos se declara, con las dos versiones del resultado. Como
trabajo complementario (sección~\ref{sec:complementario}), una reducción de dimensionalidad
muestra que mis partidas forman un continuo y no grupos, un análisis exploratorio ubica de
dónde sale el déficit de mis derrotas, y un modelo de probabilidad de ganar, elegido entre
tres familias, se publica como una página que se usa con el marcador del juego.

% ===========================================================================
\section{Descripción del problema}

Al terminar una partida, el juego muestra «DERROTA» y un marcador: KDA, daño, oro. Es
la primera evidencia que todos usamos para explicar el resultado, y es la más pobre,
porque \textbf{describe el final en lugar de explicarlo}. Una partida perdida puede
haberse decidido en el minuto 34 o en el 8, y el marcador final es el mismo.

\noindent La pregunta del proyecto es:
\begin{quote}
\emph{¿Mis partidas se pierden por un gran error final, o por la acumulación de
decisiones anteriores que nadie registra como error? ¿Y existe un minuto, medible, en
que la derrota deja de ser reversible?}
\end{quote}

\noindent\textbf{Lo que el proyecto no busca:} culpables, estado emocional
(\emph{tilt}), predicción de partidas futuras ni causalidad. El dato es observacional y
de equipo: describe las derrotas \emph{de mis equipos}, no mi aporte individual, que
no se puede separar con esta fuente. La unidad de análisis es \textbf{un jugador en
\nCuentas{} cuentas}, todas jugadas por mí en el servidor LAN.

\figura{v1_partida_ejemplo}{1}{V1 · Una derrota mía, elegida por regla y no a ojo: la
partida de A2 cuya curva \(v(t)\) está más cerca de la mediana de su arquetipo (entre
\vUnoN{} derrotas). Mi equipo fue arriba en oro \vUnoArriba{} de los \vUnoDurMin{}
minutos: abajo solo en el minuto \vUnoAbajoAntes{} y desde el \vUnoDdos. La ventaja llegó
a \vUnoMax\,\% (\vUnoMaxOro{} de oro) en el minuto \vUnoMaxMin, bajó a \vUnoAntesRemate\,\%
en el \vUnoMinAntesRemate, y en los últimos 2 minutos el rival tomó \vUnoRemateRival{}
objetivos. Arriba, los objetivos de mi equipo; abajo, los del rival ({\codigo ◆} monstruo
épico, {\codigo ■} estructura).}

% ===========================================================================
\section{Hipótesis iniciales}

En la primera entrega se plantearon tres hipótesis, reducidas a propósito desde seis:
tres pruebas al 5\,\% dan 14\,\% de probabilidad de al menos un falso positivo, y seis
daban 26\,\%. Cada una iba con la trampa que podía arruinarla.

\begin{itemize}[leftmargin=1.4em,itemsep=.4em]
  \item \textbf{H1 · Ventaja desaprovechada.} Con la misma ventaja temprana, mi equipo
        toma menos objetivos por minuto en las partidas que perdemos que en las que
        ganamos. \emph{Trampa:} «desperdiciamos la ventaja» es tautológico si se mira
        el resultado; se mide sin mirarlo.
  \item \textbf{H2 · El quiebre coincide con un objetivo.} El minuto en que más cae la
        ventaja (lo que después se llama D1) coincide con un objetivo mayor del rival
        (barón, dragón, heraldo, larvas, Atakhan) más seguido de lo que explicaría el
        azar. \emph{Trampa:} el objetivo marca el momento, no lo causa.
  \item \textbf{H3 · Punto de no retorno.} Existe un minuto a partir del cual mis
        derrotas ya casi nunca se revierten. \emph{Trampa:} sesgo de supervivencia; el
        denominador se encoge justo donde estaría el efecto.
\end{itemize}

\noindent Además se registraron cuatro \textbf{arquetipos de derrota} con umbrales
numéricos (sección~\ref{sec:arq}) y un plan de nueve experimentos (E0--E8), todo antes
de descargar mis partidas.

% ===========================================================================
\section{Bases de datos}

\begin{itemize}[leftmargin=1.4em,itemsep=.3em]
  \item \textbf{Riot Games API, \texttt{match-v5}} (fuente principal). Por partida, el
        \emph{registro} (resultado, campeones, roles, estadísticas finales) y la
        \emph{línea temporal} (un cuadro por minuto con el oro de los diez jugadores,
        más cada evento con su marca de tiempo en milisegundos).
  \item \textbf{Data Dragon} (apoyo): catálogo estático para nombres e imágenes.
  \item \textbf{Corpus piloto}: 7 partidas públicas de ejemplo, usadas en el EDA y en
        la actividad de anonimización; aquí solo como antecedente.
\end{itemize}

\noindent\textbf{Descarga.} Un script propio pidió, por cada cuenta, la lista completa
de partidas (\listaTotal{} en total) y la de cola 420 explícita (\listaCola), y
descargó registro y línea temporal de la más vieja a la más nueva, retomando si la
llave caducaba o se caía la red. Quedaron \nUnicas{} partidas únicas: 915\,MB de JSON
que nunca salen de la computadora.

\noindent\textbf{Seis cuentas.} Con la primera cuenta sola había \nCuentaUno{} partidas
(\derCuentaUno{} derrotas), por debajo del criterio de continuar fijado en el plan
(unas 150 partidas). Con las seis hay \nFinal{} partidas y \nDerrotas{} derrotas, y se
pudieron conservar los cuatro arquetipos.

\begin{table}[H]
\centering\small
\caption{La muestra por cuenta, después del filtrado.}
\label{tab:cuentas}
\begin{tabular}{lrrrr}
\toprule
Cuenta & Cola 420 en la lista & Muestra final & Victorias & Derrotas \\
\midrule
C1 & 103 & 102 & 61 & 41 \\
C2 & 552 & 529 & 289 & 240 \\
C3 & 191 & 188 & 99 & 89 \\
C4 & 166 & 156 & 84 & 72 \\
C5 & 141 & 132 & 76 & 56 \\
C6 & 98 & 88 & 49 & 39 \\
\midrule
\textbf{Total} & \textbf{1\,251} & \textbf{1\,195} & \textbf{658} & \textbf{537} \\
\bottomrule
\end{tabular}
\end{table}


\noindent\textbf{Retención, medida.} El plan suponía que la línea temporal dura un año
en la API. Ninguna de las \nUnicas{} pedidas devolvió 404, y la más vieja es de
\mesInicio: el supuesto era falso, en la dirección favorable.

\noindent\textbf{Cobertura.} De \mesInicio{} a \mesFin, \nParches{} parches del juego.
Mi rol es BOTTOM en \nBottom{} de las \nFinal{} partidas (\pctBottom\,\%).

% ===========================================================================
\section{Etapa de filtrado}

Hubo filtrado, con reglas fijadas antes de analizar y cada paso reportado:

\begin{enumerate}[leftmargin=1.6em,itemsep=.25em]
  \item Cola 420 confirmada en el registro (la lista de «ranked» mezcla la flexible).
  \item Diez participantes: una partida de 8 da una diferencia de oro falsa sin fallar.
  \item Sin \emph{remakes}, con \texttt{gameEndedInEarlySurrender}. No se filtra por
        \texttt{endOfGameResult}, que Riot marca mal en la cola 420.
  \item Duración de al menos 5 minutos.
  \item Sin partidas abandonadas: rendición antes del minuto 10. Regla agregada al
        revisar la primera descarga, antes de calcular cualquier resultado.
\end{enumerate}

\noindent Las victorias \textbf{no} se filtran: son el grupo de control.

\begin{table}[H]
\centering\small
\caption{Embudo de filtrado, rehecho desde la tabla pública de partidas.}
\label{tab:embudo}
\begin{tabular}{lrr}
\toprule
Paso & Partidas & Quita \\
\midrule
Con registro y línea temporal & 1\,241 & — \\
Cola 420 confirmada en el registro & 1\,241 & 0 \\
Diez participantes & 1\,241 & 0 \\
Sin remake & 1\,204 & 37 \\
Duración de al menos 5 minutos & 1\,202 & 2 \\
Sin rendición antes del minuto 10 & 1\,195 & 7 \\
\bottomrule
\end{tabular}
\end{table}


\figura{v8_embudo}{.9}{V8 · Embudo de filtrado: de \nUnicas{} partidas quedan
\nFinal{} (\nVictorias{} victorias y \nDerrotas{} derrotas).}

% ===========================================================================
\section{Preprocesamiento}

\subsection{Anonimización primero}
Cada partida trae los identificadores de los diez jugadores, nueve de ellos
desconocidos. Antes de que cualquier dato llegue a un CSV o a una gráfica se suprimen
nombres y Riot ID, se seudonimizan \texttt{puuid} y \texttt{matchId} con HMAC-SHA256 y
una llave efímera, y la fecha se generaliza a mes. Un barrido busca los \nPuuid{}
\texttt{puuid} y \nNombres{} nombres originales: los encuentra todos en el crudo
(control negativo) y \dato{\nFugas} en la salida. Como la composición de diez campeones
identifica cada partida, al repositorio solo van tablas a nivel equipo, y el análisis
completo se reproduce desde ellas. De los eventos solo se publica lo que el análisis lee:
torres y monstruos épicos, con su minuto y el equipo que los tomó (\nEventosPublicos{}
filas, sin posición ni víctima).

\noindent\textbf{Riesgo que queda.} Las tablas públicas no tienen nombres, pero no son
anónimas respecto a mí: el mes, la cuenta y la duración al segundo señalan una sola
partida en el \pctPartidaUnica\,\% de las \nPartidasPublicas{} filas. Quien conozca mis
cuentas podría ligar cada fila con mi historial público y, desde ahí, ver a los otros
nueve jugadores; pero ese historial ya es público, así que el repositorio no les agrega
exposición. La defensa es no publicar mis Riot ID. La duración no se redondea porque
cambiaría el embudo y la partida V1 sin proteger nada.

\subsection{Del JSON a la medida \(v(t)\)}
\begin{enumerate}[leftmargin=1.6em,itemsep=.3em]
  \item \textbf{Cuatro tablas}: partidas, participantes, cuadros y eventos; en la
        muestra final, \nCuadros{} jugador-minuto y \nEventos{} eventos.
  \item \textbf{Cuál de los diez soy yo}, decidido antes del HMAC. En \nDuo{}
        partidas aparecen dos de mis cuentas en el mismo equipo: otra persona jugaba
        una de ellas en dúo conmigo; yo soy la de BOTTOM.
  \item \textbf{Ventaja relativa}
        \(v(t) = \dfrac{\text{oro de mi equipo} - \text{oro del rival}}{\text{oro total de los diez}}\),
        porque 2\,467 de oro son el 7.7\,\% del mapa en el minuto 10 y el 2.3\,\% en el 28.
  \item \textbf{Cuadros indexados por posición}, nunca por \texttt{timestamp // 60000},
        que junta dos minutos en una casilla.
  \item \textbf{Estructuras atribuidas por \texttt{killerId}}: en
        \texttt{BUILDING\_KILL}, \texttt{teamId} nombra al equipo que \emph{perdió} la
        torre. Con la regla correcta, las torres cuadran con el conteo oficial en las
        \nFinal{} partidas.
  \item \textbf{Objetivos}: monstruos épicos (dragón, larvas, heraldo, barón, Atakhan)
        y estructuras; no las placas. Se excluyen \neutrales{} monstruos que no tomó
        ningún equipo.
\end{enumerate}

\figura{v_mediana}{1}{\(v(t)\) mediana y rango intercuartil. En la mediana, mis
derrotas no van ganando y luego se caen: quedan abajo desde los primeros minutos. Solo
\derrotasArriba{} de \nDerrotas{} terminaron arriba en oro.}

% ===========================================================================
\section{Modelos y experimentos}
\label{sec:exp}

El encuadre es una \textbf{taxonomía descriptiva pre-registrada con intervalos}, no
una batería de pruebas. Las proporciones llevan intervalos de Wilson; las tasas,
bootstrap remuestreando \emph{partidas} (10\,000 réplicas); las coincidencias, pruebas
de permutación (10\,000). Semilla 20261010 en todo.

\subsection{E1 · ¿Qué es «el minuto del quiebre»?}
Se registraron dos definiciones: \textbf{D1}, el minuto de la mayor caída de \(v\) en
un cuadro, y \textbf{D2}, el último minuto en que \(v\) cruza de \(\ge 0\) a \(< 0\).
Difieren por más de 2 minutos en el \dato{\eUnoPct\,\%} de las derrotas (IC 95\,\%:
\eUnoLo--\eUnoHi; \eUnoDesac{} de \eUnoConDdos), por \eUnoMedDif{} minutos en la mediana.

\figura{d1_minuto}{1}{D1 cae en el minuto \dUnoMedDer{} de mediana en derrotas y en
victorias por igual: al inicio hay poco oro en el mapa y una sola muerte mueve mucho
\(v\). D1 mide el ruido de los primeros minutos, no el quiebre.}

\figura{v9_d1_vs_d2}{.62}{V9 · Las dos definiciones, partida por partida. La banda
gris es el acuerdo de \(\pm2\) minutos.}

\noindent\textbf{Decisión:} D2 pasa a ser la definición principal. Ya estaba
registrada como control, así que no es una regla inventada después; D1 se reporta como
la definición que falló. La mediana de D2 en derrotas es el minuto \dDosMed. D2 es
\emph{el último minuto en que mi equipo pasa a ir abajo en oro}, y se define en
retrospectiva (hay que conocer toda la partida). No es «el minuto en que la partida
quedó perdida»: en ese momento la partida suele seguir pareja.

\subsection{E2 · H1: objetivos con ventaja}
Tasa de objetivos propios por minuto mientras hay \emph{ventaja real} (\(v \ge 5\)\,\% y
al menos 1\,500 de oro). Agregada, \hUnoDerTasa{} en derrotas contra \hUnoVicTasa{} en
victorias. Pero H1 dice «con la misma ventaja», y al ganar la ventaja es mayor y más
tardía; la comparación justa es por fase:

\begin{table}[H]
\centering\small
\caption{H1: objetivos propios por minuto con ventaja real, por fase (IC 95\,\% bootstrap por partida).}
\label{tab:h1fases}
\begin{tabular}{lrrrr}
\toprule
Fase & Derrotas & n & Victorias & n \\
\midrule
Antes del minuto 15 & 0{,}47 [0{,}40; 0{,}55] & 91 & 0{,}54 [0{,}51; 0{,}58] & 365 \\
Minutos 15 a 25 & 0{,}61 [0{,}54; 0{,}67] & 72 & 0{,}95 [0{,}92; 0{,}98] & 434 \\
Minuto 25 o después & 0{,}58 [0{,}45; 0{,}73] & 29 & 1{,}10 [1{,}05; 1{,}15] & 428 \\
\bottomrule
\end{tabular}
\end{table}


\figura{v3_h1_objetivos}{1}{V3 · Antes del minuto 15 los intervalos se traslapan; desde
el 15, las derrotas toman claramente menos objetivos con la misma ventaja.}

\subsection{E3 · H2: quiebre y objetivo del rival}
Proporción de quiebres (D2) a \(\pm1\) minuto de un monstruo épico del rival, contra una
nula que desplaza los objetivos dentro de cada partida. Observado: \hDosObsDer\,\%.

\noindent\textbf{La conclusión depende de la nula.} La registrada desplazaba los
objetivos por toda la partida; pero ningún monstruo épico existe antes del minuto
\(\approx 5\) (el primero de la muestra cae en el \hDosPrimerObj), y \hDosAntesCinco{}
quiebres ocurren antes. Contra la registrada (\hDosRegDer\,\% esperado) lo observado
queda \emph{por debajo} del azar; contra la corregida, que desplaza solo desde el minuto
5 (\hDosCorDer\,\%), queda \emph{por encima}: \dato{p = \hDosPCorDer}. En las victorias
la coincidencia es mayor todavía: \hDosObsVic\,\% contra \hDosCorVic\,\%.

\noindent\textbf{Con la definición registrada.} H2 se registró con D1, la mayor caída.
Con D1, en las derrotas lo observado es \hDosUnoObsDer\,\%: contra la nula registrada
(\hDosUnoRegDer\,\%) queda por debajo (\hDosUnoPRegDer), y contra la corregida
(\hDosUnoCorDer\,\%) no se distingue del azar (\hDosUnoPCorDer). En las victorias, otra
vez, la coincidencia es mayor (\hDosUnoObsVic\,\% contra \hDosUnoCorVic\,\%,
\hDosUnoPCorVic).

\noindent\textbf{Veredicto: no concluyente.} De las cuatro combinaciones (D1 o D2, nula
registrada o corregida), solo una apoya H2, y es la que suma los dos cambios hechos
después de ver los datos. Su \(p = \hDosPCorDer\) queda justo en el límite del umbral
corregido por tres hipótesis (\(\alpha = 0{,}05/3 = \alfaBonf\)). Lo que sí es robusto es
que la coincidencia es mayor al ganar: describe el ritmo de cualquier partida, no cómo
se pierde.

\figura{v4_h2_permutacion}{1}{V4 · Las dos nulas lado a lado, con la proporción
observada (D2). La rama con D1 se reporta en el texto.}

\subsection{E4 · H3: ¿hay un punto de no retorno?}
Para cada minuto, de las partidas que van abajo con \(v \le -5\)\,\%, cuántas terminaron
ganando, con el número en riesgo a la vista.

\begin{table}[H]
\centering\small
\caption{H3: de las partidas que van abajo con \(v \le -5\)\,\% en el minuto \(t\), cuántas terminaron ganando.}
\label{tab:h3curva}
\begin{tabular}{rrrr}
\toprule
Minuto & En riesgo & Remontaron & \% (IC 95\,\% Wilson) \\
\midrule
5 & 241 & 89 & 37 [31; 43] \\
10 & 216 & 51 & 24 [18; 30] \\
15 & 245 & 50 & 20 [16; 26] \\
20 & 216 & 26 & 12 [8; 17] \\
25 & 209 & 17 & 8 [5; 13] \\
30 & 154 & 6 & 4 [2; 8] \\
35 & 51 & 1 & 2 [0; 10] \\
\bottomrule
\end{tabular}
\end{table}


\noindent\textbf{Dónde está el punto sin regreso.} H3 dice «existe un minuto a partir del
cual mis derrotas ya casi nunca se revierten», pero el plan no fijó cuánto es «casi
nunca». Se fija aquí, y se declara como desviación: \emph{menos del \hTresUmbral\,\%}
(menos de una de cada veinte), sostenido en todos los minutos siguientes con al menos
30 partidas en riesgo. Con esa regla, el punto es el \dato{minuto \hTresPunto}: remontó el
\hTresPuntoP\,\% (\hTresPuntoRemontaron{} de \hTresPuntoRiesgo; IC 95\,\%:
\hTresPuntoLo--\hTresPuntoHi). El criterio de refutación registrado para V5 («que la curva
no baje, o que la banda sea tan ancha que no permita ubicar ningún minuto») no se cumple:
la curva baja y la banda es estrecha.

\noindent El minuto exacto depende del umbral: con 10\,\% sería el \hTresPuntoDiez, y con
3\,\% el \hTresPuntoTres{} (ya con solo \hTresPuntoTresRiesgo{} partidas en riesgo). Que se
mueva tanto es justo lo que dice la curva: antes del punto, la probabilidad de remontar
baja de forma gradual, sin un salto. Es una curva entre partidas: describe la
probabilidad de remontar en cada minuto, no la trayectoria de cada partida.

\figura{v5_h3_remontada}{.95}{V5 · La probabilidad de remontar baja de forma gradual;
desde el minuto \hTresPunto{} queda por debajo del \hTresUmbral\,\% (línea punteada).}

\subsection{E5 · Arquetipos de derrota}
\label{sec:arq}
Reglas registradas, aplicadas en orden A4 → A2 → A1 → A3 (lo que no cumple ninguna
queda sin clasificar):
\begin{itemize}[leftmargin=1.4em,itemsep=.2em]
  \item \textbf{A4 · Muy desigual:} \(v \le -10\)\,\% en algún minuto hasta el 15, y
        \(v \le -5\)\,\% desde el 15.
  \item \textbf{A2 · Ventaja desperdiciada:} \(\ge 3\) minutos seguidos de ventaja real,
        con ventaja real en el minuto 15 o después.
  \item \textbf{A1 · Derrota desde línea:} desventaja real en el minuto 15 y la
        diferencia de oro nunca vuelve a ser positiva.
  \item \textbf{A3 · Decidida en un objetivo:} \(|v| < 5\)\,\% hasta D2, y un monstruo
        épico del rival a \(\pm1\) minuto de D2.
\end{itemize}

\begin{table}[H]
\centering\small
\caption{E5: arquetipos de las 537 derrotas (IC 95\,\% de Wilson).}
\label{tab:arquetipos}
\begin{tabular}{lrr}
\toprule
Arquetipo & Derrotas & \% [IC 95\,\%] \\
\midrule
A1 · Derrota desde línea & 103 & 19 [16; 23] \\
A2 · Ventaja desperdiciada & 55 & 10 [8; 13] \\
A3 · Decidida en un objetivo & 23 & 4 [3; 6] \\
A4 · Muy desigual & 78 & 15 [12; 18] \\
Sin clasificar & 278 & 52 [48; 56] \\
\bottomrule
\end{tabular}
\end{table}


\figura{v6_arquetipos}{.95}{V6 · El residuo «sin clasificar» es la barra más alta: ese
era el criterio de refutación registrado para esta figura, y se cumplió.}

\figura{v2_arquetipos_curvas}{1}{V2 · Cada derrota por arquetipo. Las formas sí se
separan; A3 casi no existe.}

\subsection{E6 · Sensibilidad del umbral}
Con el umbral de 5\,\% en 4\,\% cambian de arquetipo \eSeisCuatro{} derrotas
(\eSeisCuatroPct\,\%); con 6\,\%, \eSeisSeis{} (\eSeisSeisPct\,\%).

\begin{table}[H]
\centering\small
\caption{E6: derrotas por arquetipo al mover el umbral de 5\,\%.}
\label{tab:sensibilidad}
\begin{tabular}{lrrrrr}
\toprule
Umbral & A1 & A2 & A3 & A4 & Sin clasificar \\
\midrule
4\,\% & 122 & 75 & 10 & 86 & 244 \\
5\,\% & 103 & 55 & 23 & 78 & 278 \\
6\,\% & 88 & 33 & 49 & 64 & 303 \\
\bottomrule
\end{tabular}
\end{table}


\figura{v7_sensibilidad}{.85}{V7 · A3 es el arquetipo más inestable.}

\subsection{E7 · Estabilidad en el tiempo}
Partiendo la muestra en \eSieteCorte{} (\eSieteUno{} derrotas antes, \eSieteDos{}
después), la mediana de D2 es el minuto \eSieteMedUno{} contra el \eSieteMedDos: no se
mueve más de 2 minutos. El reparto de arquetipos es igual en ambas mitades
(\(\chi^2\), p = \eSieteP).

\subsection{E8 · Cuentas y rol}
Las \nCuentas{} cuentas no se reparten distinto entre arquetipos (permutación,
\dato{p = \eOchoP}; la aproximación \(\chi^2\) da \eOchoPAprox, pero tiene celdas con
frecuencia esperada de \eOchoEspMin). Solo con BOTTOM (\botPartidas{} partidas,
\botDerrotas{} derrotas) todo se repite: H1 por fase \botHUnoDerA/\botHUnoDerB/\botHUnoDerC{}
contra \botHUnoVicA/\botHUnoVicB/\botHUnoVicC, H2 p = \botHDosP, H3 \botHTresTreinta\,\% en
el minuto 30.

\figura{v10_cuentas}{1}{V10 · Arquetipos por cuenta: ninguna se aparta del resto.}

% ===========================================================================
\section{Visualizaciones y recursos}
\label{sec:vis}

Cada visualización se planeó para soportar o refutar algo concreto, con su criterio de
refutación escrito antes de ver los datos. Esta tabla dice qué pasó con cada una.

\begin{longtable}{p{0.8cm}p{3.3cm}p{1.2cm}p{4.3cm}p{4.6cm}}
\caption{Las visualizaciones planeadas y su resultado.}\\
\toprule
Id & Visualización & Sirve a & Qué la haría refutar (registrado) & Qué pasó \\
\midrule
\endfirsthead
\toprule
Id & Visualización & Sirve a & Qué la haría refutar (registrado) & Qué pasó \\
\midrule
\endhead
V1 & Curva de oro de una partida, con objetivos & Problema & — (es el gancho). & Una derrota mía elegida por regla: la mediana de A2 (figura~\ref{fig:v1_partida_ejemplo}). \\
\addlinespace
V2 & \(v(t)\) de cada derrota, por arquetipo & E5 & Que los paneles se vean iguales. & Se separan; A3 casi no existe. \\
\addlinespace
V3 & Objetivos por minuto con ventaja, por fase & H1 & Intervalos que se traslapan por completo. & Se traslapan antes del minuto 15 y se separan después. \\
\addlinespace
V4 & Distribución nula de la permutación & H2 & Que la línea caiga dentro del cuerpo de la nula. & Depende de la nula y de la definición: no concluyente. \\
\addlinespace
V5 & Probabilidad de remontar por minuto & H3 & Que la curva no baje, o una banda demasiado ancha. & No se refuta: baja, y desde el minuto \hTresPunto{} queda bajo el \hTresUmbral\,\%. \\
\addlinespace
V6 & Frecuencia de arquetipos & E5 & Que «sin clasificar» sea la barra más alta. & Se cumplió: es la más alta. \\
\addlinespace
V7 & Umbral \(\times\) arquetipo & E6 & Que las partidas salten con un punto de umbral. & A3 es inestable; el resto cambia poco. \\
\addlinespace
V8 & Embudo de filtrado & Datos & — (transparencia). & De \nUnicas{} a \nFinal{} partidas. \\
\addlinespace
V9 & D1 contra D2 & E1 & Puntos lejos de la diagonal. & Se cumplió: D1 se descartó como definición principal. \\
\addlinespace
V10 & Arquetipos por cuenta & E8 & Que una cuenta se separe. & Ninguna se separa (\(p = \eOchoP\)). \\
\bottomrule
\end{longtable}

\noindent\textbf{Color.} Un solo código en todo el proyecto: azul para mi equipo y
naranja para el rival en el reporte (rojo en las presentaciones, sobre fondo oscuro), y
dorado solo para anotaciones que llevan texto. La paleta se validó con un verificador de
contraste y de separación para daltonismo antes de usarla; el signo de la ventaja
también va en la forma (arriba o abajo de la línea), no solo en el color.

\noindent\textbf{Presentaciones.} Hay dos, generadas desde los mismos resultados:
\begin{itemize}[leftmargin=1.4em,itemsep=.25em]
  \item \textbf{La presentación final} (\texttt{LoL\_PresentacionFinal}, 13 láminas,
        unos 10 minutos) es la de la exposición oral. Sigue los puntos de la rúbrica,
        con el método completo: filtrado, anonimización, las dos definiciones del
        quiebre y las dos nulas.
  \item \textbf{La presentación narrada} (\texttt{LoL\_PresentacionHistoria}, 17
        láminas) cuenta lo mismo como una historia para quien nunca ha jugado: explica
        el juego en una lámina, traduce cada gráfica a una idea cotidiana y anima cada
        gráfica paso a paso. Es material complementario, con un video de 6 minutos.
\end{itemize}
Las dos son HTML que funcionan sin internet. Las gráficas de la narrada se dibujan en el
navegador desde los datos incrustados, y las de la final son las mismas figuras del
análisis con el tema del mazo. En ninguna hay cifras escritas a mano: un script
(\texttt{construir.py}) las toma de los archivos de resultados, y falla si queda alguna
sin valor. Otro script (\texttt{shoot.py}) toma una captura de cada lámina y revisa que
nada se salga del escenario ni se encime, que cada gráfica se dibuje y que el escenario
quepa en pantallas de proyector.

\noindent\textbf{El video} de la narrada no es una grabación de pantalla: se genera
cuadro por cuadro con un reloj virtual que coloca cada animación en su instante exacto
antes de cada captura, de modo que ninguna se acelera ni se atrasa. El código de las
presentaciones y del video está en el repositorio público.

% ===========================================================================
\section{Desviaciones del plan}

Lo que cambió respecto a lo registrado antes de ver los datos, y por qué. Cuatro cambios
(la definición de E1, la nula de E3, la comparación por fase de H1 y el umbral que
ubica el punto sin regreso de H3) se decidieron \emph{después} de ver resultados. Por
eso en todos se reportan las dos versiones, la registrada y la corregida, o la
sensibilidad al umbral.

\begin{enumerate}[leftmargin=1.6em,itemsep=.25em]
  \item D2 sustituye a D1 como definición del quiebre (E1), y por consecuencia A3 usa
        D2. H2 se reporta con las dos.
  \item La nula de E3 se corrige: no puede poner objetivos donde no existen.
  \item H3 no fijaba cuánto es «casi nunca». Se fija en menos del \hTresUmbral\,\%
        sostenido, y se reportan también 3\,\% y 10\,\%.
  \item H1 se compara por fase, porque la hipótesis pide «la misma ventaja».
  \item E8 agrega un valor p por permutación, porque la aproximación \(\chi^2\) no es
        confiable con frecuencias esperadas menores a 5.
  \item Precisiones fijadas antes de cada experimento: sin placas de torre, cada larva
        es un evento, se excluyen los monstruos que no tomó nadie, y un evento se juzga
        con \(v\) del cuadro de su minuto.
\end{enumerate}

% ===========================================================================
\section{Hipótesis finales}

Con los criterios registrados en el plan (una hipótesis se mantiene si la muestra
alcanza para contrastarla, aunque el resultado sea en contra; se reformula solo por una
razón de dato), las tres se mantienen. Lo que cambia es lo que los datos dicen de cada
una.

\begin{longtable}{p{1.0cm}p{6.4cm}p{7.2cm}}
\toprule
 & Versión final & Qué cambió y por qué \\
\midrule
\endfirsthead
\toprule
 & Versión final & Qué cambió y por qué \\
\midrule
\endhead
H1 & \emph{A partir del minuto 15}, con ventaja real, mi equipo toma menos objetivos por minuto en las derrotas que en las victorias. & \textbf{Se mantiene y se precisa.} Antes del 15 no se distingue del ruido (\hUnoDerA{} contra \hUnoVicA); desde el 15, sí (\hUnoDerB{} contra \hUnoVicB). \\
\addlinespace
H2 & El quiebre (D2: el último minuto en que mi equipo pasa a ir abajo en oro, definido en retrospectiva) coincide con un monstruo del rival más de lo que explicaría el azar. & \textbf{Se mantiene; no concluyente.} Con la definición registrada (D1) no se distingue del azar (\hDosUnoPCorDer); con D2 solo supera a la nula corregida (\(p = \hDosPCorDer\), en el límite de \(\alpha = \alfaBonf\)). En las victorias coincide más. \\
\addlinespace
H3 & Existe un minuto a partir del cual mis derrotas ya casi nunca se revierten. & \textbf{Se mantiene y se confirma.} Con «casi nunca» como menos del \hTresUmbral\,\% (umbral fijado después de ver los datos), es el minuto \hTresPunto{} (\hTresPuntoP\,\%, IC 95\,\%: \hTresPuntoLo--\hTresPuntoHi). Antes, la probabilidad de remontar baja de forma gradual; por eso el minuto depende del umbral (\hTresPuntoDiez{} con 10\,\%). \\
\bottomrule
\end{longtable}

% ===========================================================================
\section{Conclusiones}

\begin{enumerate}[leftmargin=1.6em,itemsep=.45em]
  \item \textbf{Mis partidas casi nunca se pierden en el último momento: se van
        perdiendo, y hacia el minuto \hTresPunto{} ya casi no tienen regreso.} Remontar
        estando abajo pasa del \hTresCinco\,\% en el minuto 5 al \hTresQuince\,\% en el
        15, y desde el \hTresPunto{} lo logra menos de una de cada veinte partidas. Es
        una curva entre partidas: dice qué tan probable es remontar en cada minuto, no
        cómo cae cada partida.
  \item \textbf{Lo que distingue mis derrotas no es tener ventaja, sino qué hago con
        ella a mitad de partida.} Con la misma ventaja, antes del minuto 15 la tasa de
        objetivos al perder (\hUnoDerA) no se distingue de la de ganar (\hUnoVicA);
        después, las derrotas toman \hUnoDerB{} objetivos por minuto contra \hUnoVicB. Es asociación: tomar
        objetivos también produce oro.
  \item \textbf{Un monstruo del rival no explica el quiebre.} Con la definición
        registrada no hay coincidencia; con la principal es débil y depende de la nula; y
        al ganar es mayor que al perder. Describe el ritmo de cualquier partida, no cómo
        pierdo yo.
  \item \textbf{Las cuatro formas de perder explican menos de la mitad.} El
        \pctSinClas\,\% sin clasificar son, sobre todo, partidas parejas que se deciden
        tarde (observación exploratoria, no hipótesis): en la mediana pasan el
        \exTranquilo\,\% del tiempo con \(|v| < 5\)\,\% y su quiebre llega en el minuto
        \exDdos.
  \item \textbf{El método importó tanto como los datos.} Una definición registrada
        falló (D1) y una nula registrada estaba mal planteada; sin verificar supuestos,
        H2 se habría reportado al revés. Y aplicar las reglas de decisión registradas,
        en vez de la impresión de la curva, cambió el veredicto de H3.
\end{enumerate}

\subsection*{Limitaciones}
Una sola persona en \nParches{} parches; el resultado es de equipo; el dato es
observacional; las partidas se juegan en rachas y no son independientes (el bootstrap
por partida lo mitiga, no lo resuelve); y la fecha está generalizada a mes por
privacidad, así que E7 parte por mes y no por día.

% ===========================================================================
\section{Trabajo complementario}
\label{sec:complementario}

Después de cerrar las hipótesis se hicieron tres trabajos más, cada uno en su carpeta del
repositorio, con su código y su documentación. No cambian ningún resultado anterior: H1--H3
y los arquetipos se quedan como se reportaron. Sus cifras también las escribe
\texttt{cifras.py}, leídas de lo que cada trabajo generó.

\subsection{Reducción de dimensionalidad}
\textbf{Carpeta \texttt{LoL\_Dimensionalidad}} (notebook ejecutado y su PDF). Cada partida que
llegó al minuto 20 (\drN) se representó como un vector de 20 dimensiones,
\((v(1), \dots, v(20))\), y se proyectó a dos dimensiones con cuatro algoritmos: PCA, LDA,
t-SNE (con perplejidades de 5 a 150) y UMAP (con 5 a 200 vecinos). Los mapas se compararon
con la confiabilidad (\emph{trustworthiness}, que mide si los vecinos en el mapa eran vecinos
en las 20 dimensiones) y con la silueta por resultado.
\begin{itemize}[leftmargin=1.4em,itemsep=.2em]
  \item \textbf{No hay grupos separados.} En los cuatro métodos, victorias y derrotas forman
        un solo continuo ordenado por la ventaja de oro; la silueta por resultado va de
        \drSilMin{} a \drSilMax.
  \item \textbf{PCA basta.} Dos componentes guardan el \drPcaDos\,\% de la varianza, y la
        primera es la ventaja promedio de la partida (correlación \drPcUnoCorr).
  \item \textbf{El camino no agrega información.} LDA con los 20 minutos acierta el resultado
        en el \drLda\,\% de las partidas (validación cruzada), lo mismo que la regla «gana quien
        va arriba en el minuto 20» (\drVveinte\,\%).
  \item Las derrotas A2 («ventaja desperdiciada») caen entre las victorias en los cuatro
        mapas: hasta el minuto 20 parecen ganadas.
\end{itemize}

\subsection{Análisis exploratorio}
\textbf{Carpeta \texttt{LoL\_Exploratorio}.} Cuatro preguntas que no estaban registradas, así
que sus resultados son pistas, no confirmaciones. Usa las tablas por participante, solo en
local; publica agregados, y su script se detiene si alguna salida trae el identificador de
una partida o de un jugador, o el campeón de otro jugador.
\begin{enumerate}[leftmargin=1.6em,itemsep=.2em]
  \item \textbf{¿Perdí yo?} En mis derrotas jugando BOTTOM, mi duelo de línea aporta el
        \eaBotDer\,\% del déficit del equipo en el minuto 15 (\eaDefQuince{} de oro), y
        MIDDLE el \eaMidDer\,\%; en las victorias, mi línea aporta el \eaBotVic\,\% de la
        ventaja. Un minuto antes del quiebre (\eaAbajoN{} derrotas), mi duelo va abajo en
        el \eaAbajoYo\,\% de los casos y el de MIDDLE en el \eaAbajoMid\,\%.
  \item \textbf{¿Qué rodea al quiebre?} Peleas perdidas, no objetivos: hay
        \eaMuerteContra{} veces las muertes esperadas de mi equipo y \eaMuerteFavor{} las del
        rival; los dragones y las torres aparecen lo esperado, y solo destaca el barón del
        rival (\eaBaronContra). Mis muertes suben como las del equipo (\eaMiMuerte).
  \item \textbf{¿Importa el camino?} No: en ningún minuto la pérdida logarítmica cambia en más
        de \eaCaminoMax{} al agregar a la ventaja actual su cambio en los últimos 3 minutos. Una tabla
        da la probabilidad por minuto y ventaja (por ejemplo, ir \(-5\)\,\% en el minuto 15 es
        ganar \eaMenosCincoQuince{} de cada 20).
  \item \textbf{¿Dónde muero?} Desde el minuto 15, en las derrotas muero \eaMuerteDer{} veces
        por cada 10 minutos (\eaMuerteVic{} en las victorias), y solo el \eaCampoDer\,\% de esas
        muertes ocurre en campo rival (\eaCampoVic\,\% en las victorias): muero defendiendo.
\end{enumerate}

\subsection{Modelo de probabilidad de ganar}
\textbf{Carpeta \texttt{LoL\_Probabilidad}.} Una página sin internet
(\texttt{probabilidad.html}) donde se escribe el marcador como se ve en el juego (minuto,
asesinatos, torres e inhibidores, dragones y la suma de niveles de cada equipo, visible con
Tab), y opcionalmente el oro, y dice cuántas de cada 20 partidas así gané, con el rango del
modelo.

\noindent\textbf{Método.} Una fila por partida y minuto de las \pgPartidas{} partidas, y un
banco de pruebas común con reglas fijadas antes de probar nada: validación cruzada de 5
pliegues \emph{por partida} (los minutos de una partida nunca quedan repartidos entre
entrenamiento y prueba), pérdida logarítmica como medida principal (también Brier, AUC y
calibración), y una prueba temporal con las \pgPrueba{} partidas de mayo de 2026 en adelante.
Regla de elección: el modelo más simple dentro del intervalo del mejor, con intervalos por
\emph{bootstrap} remuestreando partidas. Se compararon tres familias, exploradas en paralelo:
regresión logística con efectos que cambian con el minuto, árboles combinados
(\emph{gradient boosting} y bosques aleatorios) y otras (redes neuronales pequeñas, un modelo
por minuto, vecinos cercanos).

\begin{table}[H]
\centering\small
\caption{Modelo de probabilidad de ganar: pérdida logarítmica en validación cruzada por partida (menor es mejor).}
\label{tab:probabilidad}
\begin{tabular}{lrr}
\toprule
Modelo & Solo el marcador & Marcador y oro \\
\midrule
No saber nada (la proporción de victorias) & 0{,}689 & — \\
Logística simple (con oro: solo el oro) & 0{,}554 & 0{,}545 \\
Mejor árbol combinado & 0{,}554 & 0{,}537 \\
Mejor de las otras familias & 0{,}549 & 0{,}531 \\
Logística con efectos por minuto & 0{,}548 & 0{,}531 \\
\textbf{La misma, con los niveles (elegida)} & \textbf{0{,}536} & \textbf{0{,}527} \\
\bottomrule
\end{tabular}
\end{table}


\noindent\textbf{Resultados.} Ningún árbol ni red superó a la logística por más que el margen
de error: con unas mil partidas no hay datos para aprender más. Una curva de aprendizaje
confirmó que más partidas casi no ayudan; lo que mejoró al modelo fue información nueva. La
suma de niveles bajó la pérdida de \pgAntesMarc{} a \pgElegMarc{} sin el oro (con él, de
\pgAntesOro{} a \pgElegOro), la mejora más grande de todo el proceso: con los niveles, el
marcador predice mejor que el oro solo (\pgSoloOro). La experiencia, que no se ve en la
partida, no agrega nada sobre los niveles: la pérdida cambia apenas \pgXpDif, dentro del
margen de error. El detalle de los objetivos (torres e
inhibidores por separado, alma, dragón ancestral, inhibidores caídos), leído del JSON crudo,
no mejoró nada: todas las variantes quedaron a menos de \pgObjDifMax{} del modelo elegido. El
elegido tiene \pgNumMarc{} parámetros sin el oro y \pgNumOro{} con él. Entre lo que predice y
lo que se observó hay, a lo más, \pgCalMax{}~pp de diferencia (diez grupos, fuera de muestra),
y la página calcula exactamente lo mismo que Python (lo comprueba un script en cada
construcción).

\noindent\textbf{Límites.} En la prueba temporal la pérdida subió para todos los modelos:
el elegido marcó \pgPruebaMarc{} sin el oro y \pgPruebaOro{} con él (el modelo sin niveles,
\pgPruebaAntesMarc{} sin el oro). Esa prueba se volvió a usar después de corregir los datos (una
revisión independiente encontró que las ejecuciones contaban como asesinatos) y de cada ronda,
así que ya no es una prueba ciega. El modelo describe mis partidas, no las de cualquier jugador,
y sus coeficientes son asociaciones.

% ===========================================================================
\section{Reproducibilidad}

\noindent\textbf{Repositorio público:}
\url{https://github.com/hmoran457-creator/lol-cuando-se-pierde-la-partida}, con el código, las
tablas públicas, los resultados, este documento y las presentaciones.

Todo el análisis (E1--E8 y las figuras) y este documento se reproducen desde las
\textbf{tablas públicas} anonimizadas, sin el JSON crudo ni una llave de la API: las
pocas cifras que salen de los datos locales (listas por cuenta, dúos, eventos de la
muestra) van ya contadas en \texttt{resumen\_descarga.json}. Las versiones están fijadas
en \texttt{requirements.txt} (Python 3.13.13, pandas 3.0.2, numpy 2.4.4, scipy 1.18.1,
matplotlib 3.11.1; playwright solo para revisar las presentaciones). Para el PDF hace
falta LuaLaTeX (TeX Live o MiKTeX) y la fuente DejaVu Sans Mono. Orden de ejecución,
desde la carpeta del proyecto:

\begin{Verbatim}[fontsize=\small, formatcom=\codigo, frame=lines]
# 1-4 · solo con llave propia (da otra muestra: la de quien descarga)
python LoL_Datos/scripts/descargar_partidas.py --plataforma la1   # una vez por cuenta
python LoL_Datos/scripts/revisar_descarga.py
python LoL_Anonimizacion/scripts/anonimizar.py --datos LoL_Datos/crudo
python LoL_Datos/scripts/resumen_publico.py
# 5 · desde las tablas públicas (unos 5 minutos en total)
python LoL_Analisis/scripts/e1_quiebre.py
python LoL_Analisis/scripts/e2_e4.py
python LoL_Analisis/scripts/e5_e8.py
python LoL_Analisis/scripts/v1_partida.py
# 6 · este documento
python LoL_Reporte/scripts/cifras.py
latexmk -cd -lualatex LoL_Reporte/reporte_final.tex
\end{Verbatim}

\noindent \textbf{Trabajo complementario.} Necesita además scikit-learn, umap-learn, nbclient,
nbconvert y pyarrow (versiones fijadas en el mismo \texttt{requirements.txt}). Desde el
repositorio público corren el notebook de reducción de dimensionalidad y la construcción y
revisión de la herramienta de probabilidad; el notebook exploratorio y el entrenamiento del
modelo leen las tablas por participante, que no se publican, así que se incluyen ya
ejecutados. Cada carpeta explica su orden de ejecución en su LEEME.

\begin{Verbatim}[fontsize=\small, formatcom=\codigo, frame=lines]
python LoL_Dimensionalidad/scripts/construir_notebook.py   # ejecuta el notebook y hace el PDF
python LoL_Probabilidad/scripts/construir.py               # la página, desde modelo.json
python LoL_Probabilidad/scripts/verificar.py               # paridad JS-Python y revisión
\end{Verbatim}

\noindent \textbf{Pruebas automáticas} (\texttt{python -m pytest} en
\texttt{LoL\_Datos/tests}, \texttt{LoL\_Anonimizacion/tests} y
\texttt{LoL\_Analisis/tests}): el descargador contra una API simulada, el barrido de
fugas con fugas plantadas, el clasificador de arquetipos contra casos conocidos de
antemano, y las definiciones centrales (\(v(t)\), D1, D2, el embudo y el punto sin
regreso) contra casos construidos a mano. Corren en una copia limpia del repositorio.
\textbf{Verificación independiente} (hecha durante el proyecto; no es un script del
repositorio): las cifras observadas de E1--E4 se recalcularon directamente desde el
JSON crudo, sin pasar por la anonimización, y coincidieron exactamente. Además, una
copia limpia del repositorio regeneró los resultados idénticos byte a byte.

\noindent \textbf{Código de terceros:} no se copió código de internet; las bibliotecas
usadas se citan abajo. El código y este documento se elaboraron con la asistencia de
Claude (Anthropic).

% ===========================================================================
\section*{Referencias}
\addcontentsline{toc}{section}{Referencias}
\begin{itemize}[leftmargin=1.4em,itemsep=.3em,label={}]
  \item Brier, G.~W. (1950). Verification of forecasts expressed in terms of probability.
        \emph{Monthly Weather Review}, 78(1), 1--3.
  \item Efron, B., y Tibshirani, R.~J. (1993). \emph{An Introduction to the
        Bootstrap}. Chapman \& Hall.
  \item Fisher, R.~A. (1936). The use of multiple measurements in taxonomic problems.
        \emph{Annals of Eugenics}, 7(2), 179--188.
  \item Good, P. (2005). \emph{Permutation, Parametric and Bootstrap Tests of
        Hypotheses} (3.ª ed.). Springer.
  \item Harris, C.~R., et al. (2020). Array programming with NumPy. \emph{Nature},
        585, 357--362.
  \item Hastie, T., Tibshirani, R., y Friedman, J. (2009). \emph{The Elements of
        Statistical Learning} (2.ª ed.). Springer.
  \item Hunter, J.~D. (2007). Matplotlib: A 2D graphics environment. \emph{Computing
        in Science \& Engineering}, 9(3), 90--95.
  \item Kluyver, T., et al. (2016). Jupyter Notebooks: a publishing format for
        reproducible computational workflows. En \emph{Positioning and Power in Academic
        Publishing}, 87--90. IOS Press.
  \item McInnes, L., Healy, J., y Melville, J. (2018). UMAP: Uniform Manifold
        Approximation and Projection for dimension reduction. \emph{arXiv:1802.03426}.
  \item McKinney, W. (2010). Data structures for statistical computing in Python.
        \emph{Proceedings of the 9th Python in Science Conference}, 56--61.
  \item MingweiSamuel. \emph{riotapi-schema} (ejemplos de respuestas de
        \texttt{match-v5}). \url{https://github.com/MingweiSamuel/riotapi-schema}
  \item Pearson, K. (1901). On lines and planes of closest fit to systems of points in
        space. \emph{Philosophical Magazine}, 2(11), 559--572.
  \item Pedregosa, F., et al. (2011). Scikit-learn: Machine learning in Python.
        \emph{Journal of Machine Learning Research}, 12, 2825--2830.
  \item Riot Games. \emph{Riot Developer Portal: APIs} y \emph{Data Dragon}.
        \url{https://developer.riotgames.com}
  \item Rocher, L., Hendrickx, J.~M., y de Montjoye, Y.-A. (2019). Estimating the
        success of re-identifications in incomplete datasets using generative models.
        \emph{Nature Communications}, 10, 3069.
  \item Rousseeuw, P.~J. (1987). Silhouettes: a graphical aid to the interpretation and
        validation of cluster analysis. \emph{Journal of Computational and Applied
        Mathematics}, 20, 53--65.
  \item Sweeney, L. (2002). k-anonymity: A model for protecting privacy.
        \emph{International Journal of Uncertainty, Fuzziness and Knowledge-Based
        Systems}, 10(5), 557--570.
  \item van der Maaten, L., y Hinton, G. (2008). Visualizing data using t-SNE.
        \emph{Journal of Machine Learning Research}, 9, 2579--2605.
  \item Venna, J., y Kaski, S. (2001). Neighborhood preservation in nonlinear projection
        methods: an experimental study. En \emph{Artificial Neural Networks — ICANN 2001},
        LNCS 2130, 485--491. Springer.
  \item Virtanen, P., et al. (2020). SciPy 1.0: Fundamental algorithms for scientific
        computing in Python. \emph{Nature Methods}, 17, 261--272.
  \item Wilson, E.~B. (1927). Probable inference, the law of succession, and
        statistical inference. \emph{Journal of the American Statistical
        Association}, 22(158), 209--212.
\end{itemize}

\vspace{.5cm}
\noindent{\small\color{suave} Este proyecto no está avalado por Riot Games y no refleja
las opiniones de Riot Games ni de nadie involucrado oficialmente en la producción o
administración de League of Legends. League of Legends y Riot Games son marcas
registradas de Riot Games,~Inc.}

% ===========================================================================
\newpage
\appendix
\section{Código}
Los scripts completos del análisis y de este documento, en el orden en que se ejecutan.
Se reproducen sin cambios. El código de las presentaciones y del video (los
\texttt{construir.py}, \texttt{shoot.py} y \texttt{video.py}) y el del trabajo
complementario (carpetas \texttt{LoL\_Dimensionalidad}, \texttt{LoL\_Exploratorio} y
\texttt{LoL\_Probabilidad}) están en el repositorio público.

\codigofuente{A.1 · Descarga: \texttt{descargar\_partidas.py}}{../LoL_Datos/scripts/descargar_partidas.py}
\codigofuente{A.2 · Revisión y embudo: \texttt{revisar\_descarga.py}}{../LoL_Datos/scripts/revisar_descarga.py}
\codigofuente{A.3 · Anonimización: \texttt{anonimizar.py}}{../LoL_Anonimizacion/scripts/anonimizar.py}
\codigofuente{A.3b · Resumen publicable de lo local: \texttt{resumen\_publico.py}}{../LoL_Datos/scripts/resumen_publico.py}
\codigofuente{A.4 · Medida \(v(t)\) y E1: \texttt{e1\_quiebre.py}}{../LoL_Analisis/scripts/e1_quiebre.py}
\codigofuente{A.5 · E2--E4: \texttt{e2\_e4.py}}{../LoL_Analisis/scripts/e2_e4.py}
\codigofuente{A.6 · E5--E8: \texttt{e5\_e8.py}}{../LoL_Analisis/scripts/e5_e8.py}
\codigofuente{A.7 · V1: \texttt{v1\_partida.py}}{../LoL_Analisis/scripts/v1_partida.py}
\codigofuente{A.8 · Cifras del reporte: \texttt{cifras.py}}{scripts/cifras.py}

\end{document}
